\begin{explanationbox}
  \textbf{\textcolor{CExplanation}{一句话直觉}}
  把两条互相垂直的切线看成一个直角的两条边，交点是 $P$。原点 $O$ 到这两条边的距离分别是 $d_1,d_2$。由勾股定理，$OP^2=d_1^2+d_2^2$；椭圆的特殊性质使这个和始终等于 $a^2+b^2$。

  \medskip

\begin{center}
\begin{minipage}{\linewidth}
\centering
\textbf{垂足矩形把两条垂距接到 $OP$}\par\smallskip
\begin{tikzpicture}[scale=1.1, every node/.style={font=\small}, line cap=round]
\draw[->,gray] (-4,0)--(4.2,0) node[right] {$x$};
\draw[->,gray] (0,-3.8)--(0,4) node[above] {$y$};
\draw[CStatement,thick] (0,0) ellipse [x radius=3,y radius=2];
\node[below left] at (0,0) {$O$};
\coordinate (O) at (0,0);
\coordinate (H) at (2.5805066,1.2033100);
\coordinate (K) at (-0.9348399,2.0047706);
\coordinate (P) at (1.6456667,3.2080806);
\coordinate (T) at (2.8647647,0.5937163);
\coordinate (U) at (-1.7194984,1.6388784);
\draw[CProof,thick] ($(H)!-1.15!(P)$)--(P);
\draw[CProof,thick] ($(K)!-.9!(P)$)--(P);
\fill (T) circle (1.7pt); \fill (U) circle (1.7pt);
\fill (P) circle (2pt) node[above right] {$P(x_0,y_0)$};
\draw ($(P)!.09!(H)$)--($(P)!.09!(H)+.09*(K)-.09*(P)$)--($(P)!.09!(K)$);

\draw[CExample,thick] (0,0)--(H) node[midway,below right] {$|c|$};
\draw[CExample,thick] (0,0)--(K) node[midway,below left] {$|d|$};
\draw[CExplanation,thick] (0,0)--(P);
\node[right] at (H) {$H$}; \node[above left] at (K) {$K$};
\draw pic[draw,angle radius=7pt] {right angle = P--H--O};
\draw pic[draw,angle radius=7pt] {right angle = O--K--P};
\node[fill=white] at (0,-2.8) {$OP^2=OH^2+HP^2=c^2+d^2$};

\end{tikzpicture}
\par $H,K$ 是垂足；矩形 $OHPK$ 给出 $HP=OK=|d|$。
\end{minipage}
\end{center}

  图中矩形把两条垂距与 $OP$ 连接起来；轴向特例见首页。接下来只需证明垂距的平方和不随方向变化。

  \medskip
  \textbf{\textcolor{CExplanation}{工具一：椭圆切线的一般式条件（高中判别式即可证明）}}
  对 $A,B$ 不同时为 $0$ 的直线 $Ax+By=C$，有
  \[
    \boxed{\text{直线与椭圆相切}\ \Longleftrightarrow\ C^2=a^2A^2+b^2B^2.}
  \]
  \textbf{为什么成立？}当 $B\ne0$ 时，把 $y=(C-Ax)/B$ 代入椭圆，得到
  \[
    (b^2B^2+a^2A^2)x^2-2a^2ACx+a^2(C^2-b^2B^2)=0.
  \]
  这是一元二次方程。直线与椭圆只有一个交点时，判别式为零：
  \[
    \Delta=4a^2b^2B^2\bigl(a^2A^2+b^2B^2-C^2\bigr)=0.
  \]
  因为 $a,b>0$ 且 $B\ne0$，故 $C^2=a^2A^2+b^2B^2$。当 $B=0$ 时，直线为 $x=C/A$，它与椭圆相切当且仅当 $x=\pm a$，条件仍是 $C^2=a^2A^2$。因此公式对\textbf{竖直、水平、倾斜}三种方向都适用。

  \medskip
  \textbf{\textcolor{CExplanation}{工具二：两条垂直线怎样把勾股定理带进来}}
  假设第一条切线为
  \[
    l_1:\ x\cos\theta+y\sin\theta=c,
  \]
  与它垂直并且过同一点 $P$ 的第二条直线可写成
  \[
    l_2:\ -x\sin\theta+y\cos\theta=d.
  \]
  因为 $\cos^2\theta+\sin^2\theta=1$，点 $O$ 到这两条直线的距离分别是 $|c|$ 和 $|d|$。它们互相垂直并交于 $P$，所以
  \[
    OP^2=c^2+d^2.
  \]
  两条直线都与椭圆相切时，由工具一可得
  \[
    \begin{aligned}
       c^2&=a^2\cos^2\theta+b^2\sin^2\theta,\\
       d^2&=a^2\sin^2\theta+b^2\cos^2\theta.
    \end{aligned}
  \]
  两式相加，$\sin^2\theta+\cos^2\theta=1$ 出现两次，于是
  \[
    OP^2=c^2+d^2=a^2+b^2.
  \]
  \textbf{几何意义：}单独看，原点到某一条切线的距离会随方向改变；但当另一条切线始终与它垂直时，一条垂距的平方增加，另一条相应减少，\textbf{总和不变}。

  \medskip
  \textbf{\textcolor{CExplanation}{工具三：从“垂直”转成“斜率”}}
  斜率存在的切线 $y=kx+m$，由工具一取 $A=-k$、$B=1$、$C=m$，得到
  \[
    \boxed{m^2=a^2k^2+b^2}.
  \]
  若它经过 $P(x_0,y_0)$，则 $m=y_0-kx_0$，即
  \[
    \boxed{(x_0^2-a^2)k^2-2x_0y_0k+(y_0^2-b^2)=0.}
  \]
  \textbf{几何意义：}这个方程的每个实根代表一条\emph{斜率存在}的过 $P$ 切线。两条斜率都存在时，利用韦达定理计算 $k_1k_2$，再使用垂直条件 $k_1k_2=-1$，不必分别求出切线。

  \medskip
  \textbf{\textcolor{CExplanation}{为什么必须讨论竖直切线}}
  切线 $x=\pm a$ 的斜率不存在，不能代入 $k_1k_2=-1$。它的垂线必然是水平线 $y=\pm b$，交点就是 $(\pm a,\pm b)$，仍满足 $x^2+y^2=a^2+b^2$。因此\textbf{用直线的一般式可避免遗漏；用斜率法时必须补充这一类}。

  \medskip

  \medskip
  \textbf{\textcolor{CExplanation}{工具四：从切点弦出发，只求两组乘积}}
  设 $P(x_0,y_0)$ 在椭圆外，两切点分别为 $A(x_1,y_1)$、$B(x_2,y_2)$。椭圆在切点 $A$ 处的切线方程为
  \[
    \frac{x_1x}{a^2}+\frac{y_1y}{b^2}=1;
  \]
  在切点 $B$ 处同理。两条切线都经过 $P$，故两个切点都在直线
  \[
    \boxed{\frac{x_0x}{a^2}+\frac{y_0y}{b^2}=1}
  \]
  上。这条经过两切点的直线叫\textbf{切点弦 $AB$}。\textbf{注意：}$AB$ 是连接两个切点的直线，本身通常不是切线。

  \medskip
  \textbf{为什么不必求四组和与积？}两条切线的一般式分别是
  \[
    \frac{x_1}{a^2}x+\frac{y_1}{b^2}y=1,\qquad
    \frac{x_2}{a^2}x+\frac{y_2}{b^2}y=1.
  \]
  对一般式 $A_1x+B_1y=C_1$ 与 $A_2x+B_2y=C_2$，垂直条件为
  \[
    A_1A_2+B_1B_2=0.
  \]
  因此
  \[
    \boxed{PA\perp PB\ \Longleftrightarrow\
    \frac{x_1x_2}{a^4}+\frac{y_1y_2}{b^4}=0}.
  \]
  \textbf{几何意义：}判断垂直需要的是\emph{两条切线的方向}。在切线的一般式中，$x,y$ 的系数已经确定了直线方向，所以\textbf{只需求 $x_1x_2$ 与 $y_1y_2$，不用求四组和与积}。

  \medskip
  \textbf{联立后的速用结果（详细推导见证明方法五）：}记
  \[
    D=b^2x_0^2+a^2y_0^2>0,
  \]
  则
  \[
    \boxed{\frac{x_1x_2}{a^4}=\frac{b^2-y_0^2}{D}},\qquad
    \boxed{\frac{y_1y_2}{b^4}=\frac{a^2-x_0^2}{D}}.
  \]
  把这两式相加恰好得到 $(a^2+b^2-x_0^2-y_0^2)/D$。\textbf{前面出现平方和，是两条切线相互垂直；最后出现分子为零，就是蒙日圆方程。}

  \medskip
  \textbf{两种方法对照：}原本用 $PA$、$PB$ 的斜率乘积 $-1$，要先算 $x_1+x_2$、$x_1x_2$、$y_1+y_2$、$y_1y_2$，还需要担心分母为 $0$；改用直线一般式后，只需要两个乘积，并且水平切线、竖直切线都能包含。

  \textbf{\textcolor{CExplanation}{考场识别信号}}
  \begin{itemize}[leftmargin=*, itemsep=0.3em]
    \item 题目出现“椭圆外一点所引两条切线互相垂直”，优先写 $OP^2=a^2+b^2$。
    \item 已知切点弦并要用韦达定理证明垂直时，先看最终需要的量：只求 $x_1x_2$、$y_1y_2$，不用先求两组坐标和。
    \item 题目要求逆向判断，先看 $P$ 是否满足蒙日圆方程，再说明两条切线存在且垂直。
    \item 出现一条竖直切线时，直接用 $x=\pm a$、$y=\pm b$ 验证，切勿使用不存在的斜率。
  \end{itemize}
\end{explanationbox}
